% Capitolo sulle prestazioni dell'ibridazione.
% Riscrittura del capitolo "Prestazioni e limiti dell'apprendimento": l'impianto passa da
% "verificare se l'ibridazione migliora i tempi" a "caratterizzare l'insieme delle grandezze
% che l'ibridazione modifica, e derivarne un criterio di decisione".
% Da aggiungere: condizioni al contorno, meccanismo, confronto con la letteratura, conclusioni.

\section{Prestazioni e limiti dell'ibridazione}

Il capitolo precedente ha stabilito che la correzione appresa migliora la previsione del
moto in modo misurabile. Resta da verificare ciò che davvero interessa: se il sistema, nel
suo complesso, ne tragga un beneficio osservabile. Le due domande non coincidono, perché
fra la previsione e il comportamento del robot si interpongono la costruzione della mappa
di costo, la ricerca del percorso e la dinamica del moto, e ciascuno di questi passaggi
può assorbire il guadagno prima che si traduca in azione.

Questa verifica non è un adempimento formale: è il criterio con cui si decide. Per chi
gestisce un sistema robotico già funzionante, aggiungere una componente appresa è un
investimento --- sviluppo, raccolta e mantenimento dei dati, capacità di calcolo a bordo
--- e ha senso soltanto a fronte di un ritorno misurabile. Un modello più accurato che non
modifichi il comportamento del sistema non giustifica quella spesa, per quanto accurato
sia; ed è precisamente per questo che il miglioramento della previsione, da solo, non
costituisce una risposta.


\subsection{Le variabili in gioco}

Formulare la questione come \emph{se l'ibridazione funzioni} è tuttavia mal posto, perché
presuppone che il sistema possieda un'unica dimensione di prestazione. Non è così.
L'introduzione di una componente predittiva appresa agisce simultaneamente su un insieme
di grandezze eterogenee, che non si muovono nella stessa direzione e che è necessario
rilevare separatamente.

\paragraph{Accuratezza della previsione.} È la grandezza direttamente ottimizzata
dall'addestramento, e l'unica su cui un miglioramento è atteso per costruzione. Si misura
come distanza fra la posizione prevista e quella effettivamente occupata dal pedone allo
scadere dell'orizzonte di previsione. È una grandezza \emph{interna}: descrive la qualità
di una stima intermedia, non un comportamento osservabile del robot.

\paragraph{Tempo di percorrenza.} È il tempo impiegato a completare un tragitto assegnato,
e costituisce la grandezza operativa per eccellenza: determina quanti incarichi il robot
porta a termine in un turno, e con essa il ritorno economico del sistema. A differenza
della precedente è una grandezza \emph{esterna}, che descrive ciò che il robot fa e non
ciò che il robot crede.

\paragraph{Tempo trascorso in arresto.} Scompone la precedente e distingue due modi
opposti di perdere tempo: fermarsi davanti a un ostacolo mobile, oppure aggirarlo
percorrendo una rotta più lunga. Senza questa scomposizione i due meccanismi risultano
indistinguibili, e poiché agiscono in direzioni opposte possono compensarsi lasciando il
tempo totale invariato pur avendo modificato profondamente il comportamento.

\paragraph{Tasso di completamento.} La frazione di esecuzioni in cui il robot raggiunge la
meta entro un limite prefissato. Un sistema che riduca il tempo medio ma fallisca più
spesso non costituisce un miglioramento, e la sola media dei tempi non lo rivelerebbe.

\paragraph{Margine spaziale rispetto alle persone.} La distanza mantenuta dai pedoni
durante il transito, rilevabile come minimo assoluto, come media dei minimi istantanei, o
come frazione di tempo trascorsa entro una soglia di prossimità. È l'indicatore prossemico
del comfort di chi condivide l'ambiente con il robot: un robot che sfiori le persone
compie il proprio incarico, ma degrada la qualità dello spazio in cui opera.

\paragraph{Fluidità del moto.} Misurabile attraverso le derivate superiori della
traiettoria, in particolare il \emph{jerk}, cioè la derivata terza della posizione.
Concorre alla leggibilità del comportamento: una traiettoria regolare consente ai pedoni
di prevedere le intenzioni del robot e di coordinarsi con esso, mentre un moto a scatti
induce esitazione e manovre reciproche inefficienti.

\paragraph{Costo computazionale.} La capacità di calcolo richiesta a bordo per l'inferenza
del modello e per il mantenimento delle strutture che ne consumano l'uscita. Nei sistemi
embedded questa risorsa compete con le altre funzioni del robot, e un modello che
richiedesse una frequenza di pianificazione inferiore potrebbe peggiorare le prestazioni
proprio in virtù della propria accuratezza. Questa grandezza non è quantificata nel
presente lavoro, ed è dichiarata qui fra le variabili in gioco perché la sua omissione
costituisce un limite noto dell'analisi e non una sua svista.

\bigskip

Queste grandezze sono in parte accoppiate, e l'accoppiamento non è casuale. Un robot che
mantenga margini più ampi percorre necessariamente traiettorie più lunghe; un robot che
anticipi i conflitti rallentando li evita ma paga il rallentamento in tempo; un robot che
ricalcoli il percorso con maggiore frequenza reagisce meglio ma consuma più calcolo. Il
guadagno di accuratezza, quando si propaga fino al comportamento, non si aggiunge dunque
al sistema come un beneficio netto: si distribuisce fra dimensioni che si contendono lo
stesso bilancio, e la forma di tale distribuzione dipende da come il pianificatore consuma
la previsione. La letteratura sulla navigazione in ambienti affollati documenta con
regolarità questo esito: l'integrazione di predittori appresi entro controllori predittivi
produce margini di sicurezza più ampi e movimento più regolare, a fronte di tempi di
percorrenza sensibilmente più lunghi \cite{aslam2025mpc}. Il guadagno non scompare, cambia
di natura.

Ne discende un'osservazione che orienta l'intero capitolo. Un lavoro che riporti soltanto
il miglioramento di accuratezza non consente di stabilire dove quel guadagno sia finito;
uno che riporti soltanto i tempi non consente di stabilire se sia stato speso altrove. La
caratterizzazione richiede l'insieme, ed è per questo che la misura descritta nel seguito
rileva tutte le grandezze elencate, comprese quelle sulle quali non si osserva alcuna
variazione: anche l'assenza di effetto su una dimensione è un'informazione sulla
distribuzione del guadagno.


\subsection{Il criterio di valutazione dipende dall'obiettivo}

Ne segue che la convenienza dell'ibridazione non è una proprietà del modello, ma una
relazione fra il modello e gli obiettivi assegnati al sistema. Un robot di servizio
logistico, valutato sul numero di consegne per turno, considera il tempo di percorrenza la
grandezza dominante, e un'ibridazione che lo peggiori del quindici per cento per guadagnare
margine costituisce per esso una perdita netta, quale che sia l'accuratezza raggiunta. Un
robot destinato a un ambiente sanitario o museale, dove la percezione di sicurezza da
parte delle persone e la prevedibilità del comportamento contano più della rapidità, può
trovare vantaggioso esattamente il medesimo scambio. La stessa componente, misurata sul
medesimo sistema, riceve verdetti opposti a seconda della funzione obiettivo adottata.

Questa dipendenza dall'obiettivo è la ragione per cui le metriche di errore comunemente
impiegate nella valutazione dei predittori non bastano a decidere: esse sono indipendenti
dal compito, e previsioni giudicate equivalenti possono condurre a esiti profondamente
diversi nella pianificazione a valle \cite{ivanovic2021rethinking, ivanovic2022injecting}.
Una valutazione utile alla decisione deve dunque misurare l'intero insieme delle grandezze
operative, non la sola accuratezza, e dichiarare rispetto a quale obiettivo il bilancio
viene tratto.


\subsection{Natura e finalità di questa analisi}

È opportuno dichiarare a questo punto la finalità del lavoro, poiché essa determina il
modo in cui i risultati vanno letti. L'obiettivo non è conseguire un sistema di navigazione
ottimale, né dimostrare che l'ibridazione produca un miglioramento: è caratterizzare in
modo completo l'effetto dell'aggiunta di una componente appresa a un sistema classico già
funzionante, rilevandone tanto i vantaggi quanto i costi e i limiti. Si tratta dunque di
uno studio di natura osservativa, non di un lavoro di ottimizzazione.

Da questa impostazione discendono due conseguenze. La prima riguarda il trattamento del
sistema di base: pianificatore, filtro di stima e logica di movimento sono mantenuti
invariati in tutte le misure, e nessun parametro viene ritarato per favorire l'una o
l'altra condizione. Modificarli renderebbe le differenze osservate non più attribuibili
alla sola correzione appresa, che è precisamente ciò che l'esperimento intende isolare;
l'assenza di interventi migliorativi sull'architettura non costituisce quindi una rinuncia,
ma una condizione necessaria dell'attribuibilità.

La seconda riguarda l'interpretazione degli esiti. In uno studio di ottimizzazione un
effetto nullo rappresenta un obiettivo mancato; in uno studio di caratterizzazione esso è
un risultato al pari di ogni altro, purché lo strumento di misura sia in grado di rilevare
effetti di quell'ordine di grandezza. Ne consegue che la validità delle conclusioni non
riposa sul segno dei valori ottenuti, bensì sulla dimostrazione della sensibilità del
protocollo sperimentale --- verifica alla quale è dedicata parte del paragrafo seguente ---
e sull'ampiezza degli intervalli entro cui gli effetti risultano confinati.

Il capitolo procede pertanto in tre tempi. Si definisce anzitutto un protocollo
sperimentale che consenta di attribuire alla sola correzione appresa le differenze
osservate, e se ne verifica la sensibilità invece di assumerla; si misura quindi l'effetto
dell'ibridazione su ciascuna delle grandezze elencate, distinguendo i risultati stabiliti
dalle tendenze non confermate; si esaminano infine le condizioni al contorno --- densità
della folla, portata sensoriale, struttura dell'ambiente, presenza stessa della previsione
--- per stabilire entro quali limiti le conclusioni si estendano. Ne discende un criterio
applicabile prima dell'investimento, che costituisce il contributo principale di questo
lavoro.


\subsection{Il sistema misurato}

L'ambiente di prova riproduce il piano secondo di Povo~1, edificio didattico
dell'Università di Trento, ricalcato sulla planimetria reale. La griglia di
rappresentazione conta $44 \times 73$ celle di 0,5~m di lato, corrispondenti a un'area di
$22 \times 36{,}5$~m; delle superfici così descritte 499~m\textsuperscript{2} risultano
calpestabili su 803~m\textsuperscript{2} lordi. L'ambiente è fortemente articolato:
l'82\% della superficie libera si trova entro una cella da un muro, e i corridoi hanno
larghezza tipica di 1,5~m. La scelta non è neutrale e sarà discussa nel seguito: è
l'ambiente in cui il robot trascorre la maggior frazione di tempo in arresto, e dunque
quello in cui un miglioramento della previsione dispone del massimo margine per
manifestarsi.

La popolazione di riferimento conta circa 349 individui compresenti, pari a una densità di
0,70~persone/m\textsuperscript{2}, ed è composta combinando le categorie comportamentali
descritte nel capitolo precedente --- pedoni ordinari, pedoni a velocità elevata, persone
stazionarie e gruppi coesi --- in proporzioni mantenute costanti in tutte le esecuzioni.
La densità complessiva è ancorata a un riferimento normativo anziché a una stima
soggettiva: il D.M. 26 agosto 1992 in materia di prevenzione incendi nell'edilizia
scolastica fissa in 26 persone l'affollamento massimo di un'aula, il che per le 25 aule del
piano determina un affollamento complessivo di 650 persone, ossia
1,30~persone/m\textsuperscript{2}. La densità impiegata corrisponde dunque a poco più
della metà dell'affollamento massimo ammesso, condizione severa ma non estrema. Due
ulteriori livelli, 0,51 e 0,64~persone/m\textsuperscript{2}, ottenuti riscalando la
medesima composizione, sono impiegati per verificare la dipendenza dei risultati dalla
densità.

La geometria del robot e dei pedoni resta quella di esercizio del simulatore: raggio del
robot 0,095~m, raggio dei pedoni 0,117~m, raggio di arresto di sicurezza 0,168~m, portata
del telemetro 2,04~m. Il robot pianifica su una griglia di risoluzione doppia rispetto a
quella dell'ambiente, con celle di 0,25~m, e l'orizzonte di previsione è fissato a un
secondo.


\subsubsection{Protocollo di misura}

Il tragitto assegnato è la traversata diagonale dell'ambiente, dalla cella libera più
prossima a un angolo a quella più prossima all'angolo opposto: è il percorso più lungo
disponibile, e quello lungo il quale il robot incontra il maggior numero di conflitti.
Tragitto e ambiente restano identici in ogni esecuzione; lasciarli variare introdurrebbe
nel confronto la difficoltà del percorso, sorgente di variabilità molto maggiore
dell'effetto che si intende osservare.

Ciò che cambia fra le due condizioni è unicamente il robot. Il confronto è
\emph{appaiato}: ogni ripetizione consiste in due esecuzioni che condividono il medesimo
seme del generatore pseudocasuale, per cui posizioni iniziali della folla, destinazioni
assegnate a ciascun pedone e intera sequenza di rumore del sensore coincidono. Si osserva
quindi la differenza \emph{dentro} ciascuna coppia, anziché lo scarto fra le medie di due
insiemi indipendenti, che sarebbe dominato dalla variabilità fra scene diverse. Il
controllo predittivo di velocità rimane disattivato in entrambe le condizioni: è il
secondo componente introdotto da questo lavoro e viene misurato separatamente, poiché
accendendolo qui una eventuale differenza non sarebbe più attribuibile alla sola
correzione appresa.

Il determinismo su cui l'appaiamento si fonda è stato verificato e non assunto. In due
occasioni distinte, campagne successive hanno ripetuto la condizione di riferimento con
gli stessi semi già impiegati: i tempi ottenuti coincidono con quelli in archivio in 50
ripetizioni su 50, con scarto massimo inferiore al centesimo di secondo, indipendentemente
dal processo che esegue la simulazione e dall'ordine di esecuzione.

Per ciascun confronto si riportano la differenza media fra le condizioni, il suo errore
standard, il numero di ripetizioni in cui ciascuna condizione ha prevalso e l'intervallo di
confidenza al 95\%. Il rapporto fra differenza media ed errore standard costituisce il
criterio di lettura adottato nel seguito: al di sotto dell'unità la differenza non è
distinguibile da zero, per quanto grande possa apparire la media; al di sopra di due la si
considera stabilita. Un'esecuzione in cui il robot non raggiunga la meta entro 150 secondi
equivalenti viene conteggiata come non conclusa e la coppia cui appartiene è scartata per
intero; nelle campagne qui riportate tale eventualità non si è mai verificata, e il tasso
di completamento è risultato unitario in entrambe le condizioni.

Va dichiarato un limite dell'appaiamento, specifico al contesto. I due robot muovono da
una scena identica, ma non appena scelgono percorsi anche di poco diversi scostano persone
diverse, e da quell'istante le due simulazioni divergono. Il seme comune controlla dunque
le condizioni iniziali, non l'intera traiettoria, e la riduzione di varianza che
l'appaiamento garantisce di norma risulta qui sensibilmente minore. È la ragione per cui
la dispersione fra ripetizioni resta ampia rispetto agli effetti attesi, e per cui il
numero di ripetizioni diventa il fattore che limita la sensibilità dell'esperimento.

\paragraph{Verifica di sensibilità.} Un esperimento che non rilevi alcun effetto è
informativo solo se si dimostra capace di rilevarne. A tale scopo la medesima procedura è
stata applicata a una grandezza il cui effetto è noto e atteso: la densità della folla.
Confrontando in modo appaiato le esecuzioni a 0,51 e a 0,70~persone/m\textsuperscript{2}
sui semi comuni, il tempo di percorrenza del robot di riferimento aumenta di 5,12~s con un
rapporto pari a 5,98, e il tempo trascorso in arresto di 3,88~s con un rapporto pari a
6,49. Lo stesso protocollo, sulle stesse esecuzioni, attribuisce alla correzione appresa
un vantaggio di 0,00~s con rapporto 0,00. Il banco di prova rileva dunque effetti reali
con rapporti compresi fra 3,6 e 6,5, e l'effetto della correzione con rapporti compresi
fra 0,2 e 0,6: un ordine di grandezza separa le due classi di risultati.


\subsubsection{Risultati}

\begin{table}[h]
\centering
\begin{tabular}{lccccc}
\hline
densità & ripetizioni & senza corr. & con corr. & differenza & rapporto \\
\hline
0,51 pers/m\textsuperscript{2} & 50 & 47,99 s & 48,15 s & $-0{,}16$ s & 0,38 \\
0,64 pers/m\textsuperscript{2} & 30 & 52,26 s & 51,58 s & $+0{,}68$ s & 0,63 \\
0,70 pers/m\textsuperscript{2} & 50 & 53,11 s & 53,27 s & $-0{,}16$ s & 0,22 \\
\hline
\end{tabular}
\caption{Tempo di percorrenza, confronto appaiato su tre densità. Valori positivi nella
colonna delle differenze indicano un vantaggio della correzione. Stima aggregata sulle 130
ripetizioni: $-0{,}14\%$, intervallo di confidenza al 95\% da $-1{,}46\%$ a $+1{,}17\%$.}
\label{tab:tempi-densita}
\end{table}

Il tempo di percorrenza non risulta modificato dalla correzione a nessuna delle tre
densità esaminate. I rapporti fra differenza media ed errore standard, compresi fra 0,22 e
0,63, collocano tutti i valori al di sotto della soglia di distinguibilità, e le
ripetizioni vinte da ciascuna condizione si ripartiscono in misura sostanzialmente
paritaria: 27 su 50, 17 su 30 e 24 su 50. L'intervallo di confidenza aggregato esclude
tanto un peggioramento superiore all'1,5\% quanto un miglioramento superiore all'1,2\%.

Il valore isolatamente favorevole della densità intermedia merita una precisazione, poiché
è l'unico a presentare segno positivo: esso è prodotto da una singola ripetizione anomala,
rimossa la quale la differenza media si riduce a $-0{,}09$~s. Non costituisce dunque
indizio di un regime di densità in cui la correzione risulti vantaggiosa.

Anche il tempo trascorso in arresto risulta invariato. Alla densità maggiore esso ammonta
al 15,8\% del tempo totale, e rappresenta pertanto il margine massimo teoricamente
recuperabile da una previsione perfetta; la correzione non ne estrae alcuna frazione
misurabile. Si osservi che, passando da 0,51 a 0,70~persone/m\textsuperscript{2}, tale
margine cresce del 68\% --- dal 9,4\% al 15,8\% --- senza che il vantaggio della
correzione ne risulti alterato: l'assenza di effetto non è dunque imputabile a una
mancanza di margine su cui agire.

\begin{table}[h]
\centering
\begin{tabular}{lcc}
\hline
grandezza & effetto della correzione & stato \\
\hline
accuratezza della previsione & $-12{,}5\%$ sull'errore & stabilito \\
tempo di percorrenza & nessuna variazione & stabilito \\
tempo trascorso in arresto & nessuna variazione & stabilito \\
tasso di completamento & invariato (unitario) & stabilito \\
margine spaziale medio & $+5{,}2\%$ ($\approx 3{,}6$ cm) & tendenza \\
esposizione alla prossimità & non confermato in replica & non stabilito \\
fluidità del moto & nessuna variazione & stabilito \\
costo computazionale & non quantificato & --- \\
\hline
\end{tabular}
\caption{Quadro complessivo delle grandezze rilevate.}
\label{tab:quadro-grandezze}
\end{table}

Il margine spaziale mantenuto rispetto ai pedoni in movimento costituisce l'unica
grandezza comportamentale sulla quale si osservi una variazione di segno definito. Su 80
ripetizioni appaiate, ripartite in due campagne condotte su semi disgiunti, la distanza
media dai corpi mobili passa da 0,689 a 0,724~m nella prima e da 0,736 a 0,774~m nella
seconda, con incremento del 5,2\% in entrambe e rapporto aggregato pari a 1,98. La
concordanza fra due campioni indipendenti conferisce a questo valore una credibilità
maggiore di quanto il solo rapporto suggerisca; esso resta tuttavia al limite della soglia
adottata, e viene pertanto riportato come tendenza e non come risultato stabilito.

Un'ulteriore grandezza di prossimità --- la frazione di tempo trascorsa entro mezzo metro
da un pedone in movimento --- aveva mostrato nella prima campagna una riduzione del 6,5\%
con rapporto 2,35. L'ipotesi è stata sottoposta a replica su 50 semi disgiunti dai
precedenti, fissandola in anticipo come unica ipotesi da verificare: la riduzione si è
ridotta al 2,9\% con rapporto 1,07. L'effetto non è dunque confermato e non viene
riportato fra i risultati. Se ne dà conto perché l'episodio illustra il rischio insito
nell'osservazione simultanea di più grandezze, e perché la replica su campioni disgiunti
costituisce la sola difesa disponibile contro di esso.

La fluidità del moto, misurata come modulo medio del \emph{jerk}, non risulta influenzata
dalla correzione: la variazione resta entro l'1\% in entrambe le campagne, con rapporti
inferiori a 0,4. Questa grandezza va peraltro interpretata con cautela nel presente
contesto, poiché il modello cinematico adottato non prevede inerzia né limiti di
accelerazione e il robot cambia direzione istantaneamente ai vertici del percorso; il
valore rilevato riflette pertanto in misura preponderante la discretizzazione della
griglia di pianificazione, e non un comportamento dinamico realistico.
